% Annexe de la partie « Revenus du capital » (refonte du 16/11/2026). Chiffres : ../pfu/sorties/calculs_pfu.txt
% (programme ../pfu/calculs_pfu.py) ; données : ../pfu/donnees/ ; vérifications : ../revue/hypotheses_pfu.md.
% Coupe du 17/11/2026 : reçoit du corps les figures pfu_n02, n04 et n09, le détail de la charge du 1 % et des
% variantes du report d'apport-cession.

Chiffres produits par \texttt{pfu/calculs\_pfu.py} (DGFiP, Insee, Banque de France, Tax Foundation). L'assiette de
2025 applique aux cases de la déclaration de 2024 l'évolution des revenus déclarés mesurée par la DGFiP ($-$12,9~\%
pour les dividendes et les intérêts, $-$26,4~\% pour les plus-values).

\subsection{Assiette, recettes et taux cumulés}

\begin{table}[!htbp]\centering\scriptsize
\caption{Assiette et recettes des prélèvements sur les revenus du capital (Md€)}\label{tab:pfu_a_assiette}
\begin{tabularx}{\linewidth}{@{}>{\raggedright\arraybackslash}X>{\raggedleft\arraybackslash}p{1.5cm}>{\raggedright\arraybackslash}p{6.2cm}@{}}
\toprule
Poste & Montant & Précisions \\
\midrule
Dividendes (case 2DC), revenus de 2024 & 50,1 & \\
Intérêts et autres distributions (2TR, 2TS) & 11,8 & \\
Plus-values et actifs numériques, moins les pertes (3VG, 3AN, 3VH) & 17,7 & \\
\textbf{Assiette du PFU, revenus de 2024} & \textbf{79,6} & y compris les 2,54~M de foyers au barème (majorant
d'écart faible) \\
\textbf{Assiette du PFU, revenus de 2025 (estimée)} & \textbf{67,0} & dividendes 43,7, intérêts 10,3, plus-values
13,0 ; contrôles : 65,5 et 72,5 \\
Plus-values placées en report d'apport-cession (3WH) & 22,2 & 33,6 en 2023 \\
Stock de reports (8UT) & 173,3 & majorant ; 127 en 2023 (évaluation préalable du PLF 2027) \\
Plus-values en sursis d'\emph{exit tax} (3WA) & 4,7 & 358 foyers \\
\midrule
IR au PFU (impôt « hors barème »), revenus de 2025 & 8,4 & 10,2 sur les revenus de 2024 (acompte de 2024 : 7,7) \\
Prélèvements sociaux sur l'assiette du PFU de 2025 & 12,5 & au taux de 2026 ; 11,5 au taux de 2025 \\
Prélèvements sociaux sur tous les revenus du capital, encaissés en 2025 & 36,2 & au taux de 17,2~\% (assiette de
210~Md€) \\
CEHR (tous revenus), revenus de 2025 & 1,4 & 1,9 sur ceux de 2024 ; part payée sur le capital inconnue \\
CDHR, revenus de 2025 & 0,62 & environ 14\,000 foyers ; acompte de 2025 : 0,37 (7\,203 foyers) \\
\bottomrule
\end{tabularx}
\par\vspace{3pt}\parbox{\linewidth}{\raggedright\scriptsize Sources : DGFiP (totaux 2042 ; Statistiques n°~41, 43
et 53) ; Assemblée nationale, rapports n°~2847 et 3056.}
\end{table}

Taux cumulés sur 100~€ de bénéfice (droit de 2026) : 48,6~\% au PFU avec l'IS de 25~\% ; avec l'IS de 25,8~\%
(contribution sociale de 3,3~\% sur l'impôt qui dépasse 763\,000~€), 49,1~\%, 52,1~\% avec la CEHR et 54,5~\% au
plancher de la CDHR ; avec l'IS de 36,1~\%, 56,2, 58,7 et 60,8~\%. Holding : 26,8~\%, 26,0~\% en intégration
fiscale. Salaire : 55,4~\% hors cotisations, comparaison favorable au capital (le bénéfice répond à un coût du travail
cotisations patronales comprises). La Tax Foundation retient pour la France l'IS
de 36,1~\% et un taux personnel de 34~\%, soit 57,9~\% ; avec l'IS de 25,8~\%, 51,0~\% en 2025 (5\textsuperscript{e}
sur 31) et 52,1~\% en 2026 (4\textsuperscript{e}, face aux taux 2025 des autres pays).

\figbloc{pfu_n04_europe}{Avec l'IS de droit commun, la France reste parmi les cinq pays européens qui taxent le plus le bénéfice distribué}{fig:pfu_europe}

\figbloc{pfu_n02_concentration}{Quatre mille foyers reçoivent un tiers des dividendes et 70~\% des plus-values imposées}{fig:pfu_concentration}

\subsection{Changer le taux : méthode et résultats}

\emph{Méthode.} L'assiette $B$ réagit de $\Delta Q = e\,B\,\Delta\ln(1-t)$, $e$ étant l'élasticité des revenus
déclarés au taux net (0,77 chez Lefebvre, Lehmann et Sicsic ; 0,65 au CAE ; 1 en hypothèse haute). Premières années :
toute la réaction compte comme recette, $t_1(B+\Delta Q) - t_0 B$ (ce qui redonne le chiffrage du CAE : 1,16~Md€
recalculé, 1,2 publié). Régime permanent : seule rapporte la part $\pi$ de la réaction qui aurait sinon échappé pour
toujours à l'impôt (purge au décès, holding jamais vidée), d'où $(t_1-t_0)B + \pi\,t_1\,\Delta Q$ ; $\pi$ n'est
mesuré nulle part.

\begin{table}[!htbp]\centering\scriptsize
\caption{Recette (+) ou coût (--) d'un changement du taux du PFU (Md€ par an)}\label{tab:pfu_a_taux}
\begin{tabularx}{\linewidth}{@{}>{\raggedright\arraybackslash}X>{\raggedleft\arraybackslash}p{1.4cm}>{\raggedleft\arraybackslash}p{3.2cm}>{\raggedleft\arraybackslash}p{3.3cm}@{}}
\toprule
Réforme, base & Mécanique & Premières années ($e$ = 0,65 / 0,77 / 1) & Permanent ($e$ = 0,77 ; $\pi$ = 0 / 0,5 / 1) \\
\midrule
PFU à 25~\%, revenus de 2025 & $-$4,29 & $-$3,32 / $-$3,14 / $-$2,79 & $-$4,29 / $-$3,71 / $-$3,14 \\
PFU à 25~\%, revenus de 2024 & $-$5,10 & $-$3,94 / $-$3,73 / $-$3,32 & $-$5,10 / $-$4,41 / $-$3,73 \\
PFU à 33~\%, revenus de 2025 & 1,07 & 0,73 / 0,67 / 0,55 & 1,07 / 0,87 / 0,67 \\
PFU à 33~\%, revenus de 2024 & 1,27 & 0,87 / 0,80 / 0,65 & 1,27 / 1,04 / 0,80 \\
\bottomrule
\end{tabularx}
\end{table}

\emph{Taux qui maximise la recette}, sans réaction croisée : $1/(1+e)$, soit 60,6 / 56,5 / 50,0~\% pour $e$ = 0,65 /
0,77 / 1 ; lever un euro de plus à 31,4~\% demande d'alourdir de 1,42 / 1,54 / 1,84~€ l'impôt calculé à comportements
inchangés (ce ne sont pas des pertes sèches). Avec la réaction des revenus du travail : 43~\% (Lefebvre, Lehmann et
Sicsic ; 57~\% sans elle).

\emph{Perte sèche d'un PFU à 25~\%.} Formule de Harberger appliquée au seul volet réel : la part $\pi$ de la réaction
compte comme perte sèche évitée, au coin moyen ; la recette perdue en régime permanent est levée ailleurs au « prix
d'un milliard » $\mu-1$ = 0,04 / 0,08 / 0,17 (barème fusionné, \texttt{synthese/sorties/architectures\_ircsg.txt}).
Avec $e$ = 0,77, revenus de 2025 : +0,17 / +0,34 / +0,73~Md€ si $\pi = 0$ ; $-$0,50 / $-$0,35 / $-$0,02 si
$\pi = 0{,}5$ ; $-$1,17 / $-$1,05 / $-$0,76 si $\pi = 1$. La baisse allège la perte sèche si la moitié au moins de
la réaction est une évasion définitive.

\emph{Bloc épargne.} Avec une élasticité réelle de l'épargne au rendement net $\varepsilon_s$, la perte sèche sur
l'assiette du PFU vaut $\tfrac12\,\varepsilon_s\,t^2/(1-t)\,B$ ; passer de 31,4 à 25~\% la réduirait de 0,20 / 0,40 /
1,01~Md€ pour $\varepsilon_s$ = 0,1 / 0,2 / 0,5 (revenus de 2025). Aucune de ces valeurs n'est estimée sur données
françaises.

\emph{Retour au barème.} Hypothèses fortes : 96~\% des dividendes au taux marginal de 45~\% et 4~\% à 30~\%, intérêts
à 30~\%, plus-values à 45~\%, abattement de 40~\% sur les dividendes. Revenus de 2025 : 10,1~Md€ en statique, avant
les foyers déjà au barème, la CDHR perdue (0,4 à 0,6) et les foyers à 41~\% ; 5,1 / 4,2 / 2,4~Md€ après réaction
($e$ = 0,65 / 0,77 / 1) ; 2,8 si les dividendes déclarés baissaient de 40~\%. L'Institut Montaigne, cité par le CAE
(2025), donne 3,2 à 3,6~Md€ en statique. Un chiffrage propre demande les fichiers POTE.

\emph{CDHR.} Rentier qui ne vit que de dividendes, à 400\,000~€ comme à 3~M€ : 38,6~\% aujourd'hui, avec un PFU à
25~\% obtenu par la part d'IR et sans la majoration de la CEHR (la CDHR reprend l'écart) ; 32,2~\% si le PFU baisse
par les prélèvements sociaux ; au barème, 38,6 et 45,3~\%. Baisser la part d'IR de 12,8 à 6,4~\% multiplie la CDHR payée par 2,05 et 2,77 : la reprise serait de l'ordre de 0,7 à 1,1~Md€ si ses
redevables étaient des rentiers purs (1,05 à 1,77 fois son rendement de 0,62~Md€), base non publiée.

\emph{Exil.} Avec 0,02 à 0,23~\% du centile aux plus hauts revenus du capital qui part à long terme par point
d'imposition effective (Bach et al., 2025, CAE), un PFU à 25~\% en retiendrait 0,13 à 1,5~\%, soit 0,02 à 0,17~Md€
de recette gardée (assiette du PFU du centile : 47~Md€ en 2025).

\subsection{La capitalisation dans une holding}

100~€ de bénéfice d'une filiale avant IS. A : distribués au PFU. B : remontés en holding, placés $N$ ans à $R$
(revenus imposés à 1,25~\%), puis distribués. C : comme B, puis transmis, les héritiers sortant la valeur sans impôt
(avec la proposition : au PFU, dix ans après). D : distribués, puis placés personnellement, revenus au PFU chaque année.
Actualisation au taux $R$.

\begin{table}[!htbp]\centering\scriptsize
\caption{Impôt total en valeur actualisée selon le chemin du bénéfice (\%)}\label{tab:pfu_a_holding}
\begin{tabularx}{\linewidth}{@{}>{\raggedright\arraybackslash}X*{6}{>{\raggedleft\arraybackslash}p{1.45cm}}@{}}
\toprule
$R$, $N$ & A & B & C & C proposition & D & Richesse B/D \\
\midrule
3~\%, 25 ans & 49,1 & 49,6 & 26,5 & 49,8 & 59,6 & +25~\% \\
3~\%, 40 ans & 49,1 & 49,9 & 26,9 & 50,0 & 64,8 & +42~\% \\
5~\%, 25 ans & 49,1 & 49,9 & 26,9 & 50,2 & 65,1 & +44~\% \\
5~\%, 40 ans & 49,1 & 50,3 & 27,6 & 50,6 & 72,1 & +78~\% \\
7~\%, 25 ans & 49,1 & 50,1 & 27,3 & 50,6 & 69,7 & +65~\% \\
\bottomrule
\end{tabularx}
\par\vspace{3pt}\parbox{\linewidth}{\raggedright\scriptsize Note : si la holding place en titres de portefeuille,
dont les revenus paient l'IS de 25~\%, B vaut 62,3~\% et C 45,0~\% ($R$ = 5~\%, $N$ = 25) : l'avantage de
capitalisation disparaît.}
\end{table}

\subsection{L'effacement à la transmission : chiffrage}

\emph{Méthode.} Plus-values effacées par an : $G = E\,b\,\lambda$ (titres concernés, part transmise chaque année,
part de plus-value latente), plus les reports transmis $S/T$ (stock sur délai moyen avant transmission). Impôt
effacé : $L = 31{,}4~\%\times G + 30~\%\times S/T$ (taux d'expiration du report non tranché). Une cohorte paie une
part $a$ de son impôt latent dans l'année, puis $(1-a)\,h\,(1-h)^{k-2}$ la $k$-ième année ($a$ = 0,1 pour les
reports) ; valeur actualisée : $L\,[a/(1+r) + (1-a)\,h/((1+r)(r+h))]$. Valeur des titres supposée constante après la
transmission (prudent).

\begin{table}[!htbp]\centering\scriptsize
\caption{Effacement à la transmission : hypothèses et résultats}\label{tab:pfu_a_effacement}
\begin{tabularx}{\linewidth}{@{}>{\raggedright\arraybackslash}X*{3}{>{\raggedleft\arraybackslash}p{2.55cm}}@{}}
\toprule
 & Bas & Central & Haut \\
\midrule
Titres concernés $E$ (Md€) & 1\,200 & 1\,450 & 1\,700 \\
Part transmise par an $b$ & 1,7~\% & 2,3~\% & 3,0~\% \\
Part de plus-value latente $\lambda$ & 40~\% & 55~\% & 70~\% \\
Ventes dans l'année $a$ ; ensuite $h$ par an & 15~\% ; 2~\% & 30~\% ; 4~\% & 45~\% ; 6~\% \\
Stock de reports $S$ (Md€) ; délai $T$ (ans) & 127 ; 35 & 144 ; 25 & 173 ; 15 \\
\midrule
Plus-values effacées par an, hors reports (Md€) & 8,2 & 18,3 & 35,7 \\
Reports transmis par an (Md€) & 3,6 & 5,8 & 11,6 \\
Impôt effacé par an (Md€) & 3,65 & 7,49 & 14,68 \\
\quad en \% du PIB & 0,12 & 0,25 & 0,49 \\
Recette, an 1 / 5 / 10 (Md€) & 0,49 / 0,74 / 1,02 & 1,90 / 2,74 / 3,62 & 5,39 / 7,43 / 9,36 \\
Recette, an 20 / 30 (Md€) & 1,50 / 1,89 & 4,91 / 5,78 & 11,81 / 13,13 \\
Valeur actualisée d'une année de transmissions, $r$ = 2,5 / 3,5 / 5~\% & 1,85 / 1,59 / 1,33 & 5,21 / 4,71 / 4,17
& 11,65 / 10,87 / 9,96 \\
\quad avec intérêt sur l'impôt différé & 3,65 & 7,49 & 14,68 \\
Clause de droits acquis : recette an 10 ; cumul ans 1-10 & 0,47 ; 2,4 & 1,43 ; 7,5 & 3,04 ; 16,4 \\
Sans clause : cumul ans 1-10 & 7,6 & 28,1 & 75,4 \\
Charge du 1~\% le plus aisé : moyenne ans 1-10 ; an 10 & 0,69 ; 0,92 & 2,53 ; 3,26 & 6,78 ; 8,42 \\
\bottomrule
\end{tabularx}
\par\vspace{3pt}\parbox{\linewidth}{\raggedright\scriptsize Note : jeux bas et haut = toutes les hypothèses basses
ou toutes les hautes (des bornes). $E$ part des 1\,969~Md€ d'actions et de fonds en actions des ménages fin 2025
(Banque de France), moins les parts de sociétés immobilières, le PEA et l'épargne salariale. Repère de $b$ : 300~Md€
de transmissions par an (CAE, 2021) pour 15\,300~Md€ de patrimoine net (Insee). Droits acquis : valeurs en hausse de
4~\% par an, reports existants exclus.}
\end{table}

\emph{Repères.} Le Trésor américain évalue l'effacement au décès à 37,8~Md\$ en 2025 (\emph{Tax Expenditures}
FY2026), soit 5,0~Md€ transposés au PIB et au taux français, sans les donations ; la Tax Foundation (Li, 2020) chiffre
sa suppression à 11 à 12~Md\$ par an, réactions comprises, soit 2,0~Md€. Le premier repère encadre l'impôt effacé, le
second la recette des premières années.

\figbloc{pfu_n10_sensibilite}{La recette de la dixième année dépend d'abord de la part transmise, de la plus-value latente et du rythme des ventes}{fig:pfu_sensibilite}

\emph{Variantes.} Une franchise de 100\,000~€ par transmission coûterait 0,4 à 1,1~Md€ par an. Déduire l'impôt latent
de l'assiette des droits de succession, exactement neutre entre vendre avant de mourir et transmettre, coûterait 0,7
à 3,0~Md€ par an selon le taux de droits effectivement payé, pacte Dutreil compris.

\subsection{La charge du 1~\% le plus aisé}

La fin de l'effacement coûterait au 1~\% 3,3~Md€ la dixième année, et 0,7~Md€ par an en moyenne avec la clause de
droits acquis. Le bilan d'ensemble (section~\ref{sec:ens_haut})
compte aussi la suppression de la majoration ex-CEHR sur le capital (0,7 à 1,1~Md€ rendus au haut) ; selon la façon
dont on répartit les revenus au PFU, la fusion peut alléger de 2,8~Md€ la charge du 1~\% ; et la fin de la
contribution exceptionnelle sur l'IS rend 5 à 7~Md€ aux quelque 450 entreprises qui la paient, dont une part
inconnue revient aux mêmes ménages.

\subsection{Le stock de reports d'apport-cession}

Valeur actualisée au taux $r$ de l'impôt sur le stock, avec trois taux de sortie annuels : expiration du vivant $e$
(environ 0,9~\% du stock en 2024), transmission $m = 1/T$, vente par les héritiers $h$. Droit actuel :
$I\,e/(r+e+m)$ ; expiration à la transmission (PLF 2027, art.~5), ou report transmis avec intérêt au taux $r$ :
$I\,(e+m)/(r+e+m)$ ; report transmis : $I\,[e+m\,h/(r+h)]/(r+e+m)$.

\begin{table}[!htbp]\centering\scriptsize
\caption{Impôt sur le stock de reports d'apport-cession, en valeur actualisée ($r$ = 3,5~\%, Md€)}\label{tab:pfu_a_reports}
\begin{tabularx}{\linewidth}{@{}>{\raggedright\arraybackslash}X*{3}{>{\raggedleft\arraybackslash}p{2.1cm}}@{}}
\toprule
 & Bas & Central & Haut \\
\midrule
Impôt nominal (30~\% du stock) & 38,1 & 43,2 & 52,0 \\
Droit actuel (effacement à la transmission) & 4,6 & 4,5 & 4,1 \\
Expiration à toute transmission (PLF 2027, art.~5) & 19,7 & 25,1 & 35,5 \\
Report transmis sans effacement (proposition) & 10,1 & 15,5 & 23,9 \\
Report transmis, avec intérêt sur l'impôt différé & 19,7 & 25,1 & 35,5 \\
Recette annuelle de l'art.~5 en régime de croisière (Md€) & 1,09 & 1,73 & 3,47 \\
\bottomrule
\end{tabularx}
\par\vspace{3pt}\parbox{\linewidth}{\raggedright\scriptsize Note : l'évaluation préalable de l'article~5 (0,5~Md€ en
2027, « en tenant compte de l'étalement ») est compatible avec une année pleine de 2 à 2,5~Md€. Borne de très long
terme : 24,7~Md€ d'entrées par an (2022-2024) $\times$ 30~\% = 7,4~Md€.}
\end{table}

\emph{L'entrepreneur qui cède} (7,5~M€ de plus-value, $h$ = 4~\% ; figure~\ref{fig:pfu_entrepreneur}) : impôt
actualisé à 3,5~\%, article~5 / report transmis, pour un décès 15, 20 ou 30 ans après : 1,41 / 0,75, 1,18 / 0,63,
0,84 / 0,45~M€ ; de 1,63 / 1,00 à 0,54 / 0,24~M€ avec un taux de 2,5 ou 5~\%. Faire expirer le report quand la
holding vend ses actifs ou devient une simple société de placement serait plus ciblé, mais difficile à contrôler et
non chiffrable. Les reports déjà transmis par donation, dont le délai de
conservation court encore, garderaient l'ancienne règle.

\figbloc{pfu_n09_entrepreneur}{Pour un entrepreneur qui cède, le report transmis vaut la moitié de l'impôt dû au décès, l'intérêt sur l'impôt différé le rétablit}{fig:pfu_entrepreneur}

\subsection{Expériences étrangères}

D'autres pays reportent la plus-value sur l'héritier (la moitié des pays de l'OCDE selon une étude citée par le
rapport n°~3056, non lue). L'évaluation des titres à la date de la réforme a compté parmi les difficultés américaines
de 1976 (non vérifié).

\begin{table}[!htbp]\centering\scriptsize
\caption{Le sort de la plus-value latente à la donation et au décès}\label{tab:pfu_a_etranger}
\begin{tabularx}{\linewidth}{@{}>{\raggedright\arraybackslash}p{1.8cm}>{\raggedright\arraybackslash}X>{\raggedright\arraybackslash}p{4.2cm}@{}}
\toprule
Pays & Règle & Source, statut \\
\midrule
France & Prix d'acquisition = valeur retenue pour les droits de mutation (CGI, art.~150-0~D) ; report effacé au
décès, et après donation si le donataire conserve six ou onze ans & Rapport n°~3056 ; évaluation préalable du PLF
2027 (lus) \\
Canada & Depuis 1972 (et la fin de l'impôt fédéral sur les successions), vente réputée faite juste avant le décès ;
conjoint au prix d'acquisition ; plus-values antérieures exonérées (« jour V ») & Agence du revenu du Canada ;
Vaillancourt et Kerkhoff (2019) (lus) ; jour V lu par résumé \\
États-Unis & Donataire : prix du donateur (IRC §~1015) ; héritier : valeur au décès (§~1014). Reprise du prix au décès
votée en 1976, abrogée en 1980. Proposition de 2015 : imposition à la donation et au décès, franchise de
100\,000~\$, report pour les entreprises familiales, paiement sur quinze ans &
\emph{U.S. Code} §~1015 ; Trésor, \emph{General Explanations} FY2016, p.~156-157 (lus) \\
Allemagne & Participations d'au moins 1~\% reçues à titre gratuit : prix du dernier acquéreur à titre onéreux (EStG,
§~17) & Loi (lue) ; autres titres non vérifiés \\
Australie & Biens acquis depuis le 20 septembre 1985 : l'héritier reprend le prix du défunt & Administration fiscale,
lue par résumé \\
\bottomrule
\end{tabularx}
\end{table}

\subsection{Hypothèses fortes et chiffres non vérifiés}

\begin{table}[!htbp]\centering\scriptsize
\caption{Hypothèses fortes (registre \texttt{revue/hypotheses\_pfu.md})}\label{tab:pfu_a_hyp}
\begin{tabularx}{\linewidth}{@{}>{\raggedright\arraybackslash}p{4.4cm}>{\raggedright\arraybackslash}p{3.3cm}>{\raggedright\arraybackslash}X@{}}
\toprule
Hypothèse & Où elle joue & Si elle est fausse \\
\midrule
Élasticité des revenus déclarés de 0,77 (0,65 à 1) & Changement de taux & À 0,5, un PFU à 25~\% coûte 3,5~Md€ les
premières années ; à 1,2, 2,5. Mesurée sur des foyers qui ont aussi des revenus d'activité \\
Réaction des dividendes surtout temporelle ($\pi$ faible) & Coût permanent d'une baisse ; perte sèche & Si
$\pi = 1$, coût permanent de 3,1~Md€ et perte sèche réduite par la baisse ; si $\pi = 0$, une hausse rapporte
durablement presque son rendement mécanique \\
Épargne insensible au rendement & « Une baisse n'achète rien » & Élasticité de 0,2 à 0,5 : la baisse gagne 0,4 à
1,0~Md€ de perte sèche, autant ou plus que son financement \\
Assiette de 2025 estimée (67~Md€) et représentative & Tous les chiffrages de taux & Celle de 2024, gonflée par la
CDHR, donne des chiffres plus élevés d'un cinquième ; celle de 2025 est peut-être déprimée par la CDHR elle-même \\
Réaction au taux attendu & Stabilité du taux & Documentée en 2024-2025 ; « suggérée » pour 2013 \\
Le bénéfice gardé en holding finit par sortir & « Report, pas exonération » & Sinon, 27~\% pour cette part ; aucune
donnée \\
$E$ = 1\,450~Md€ (1\,200 à 1\,700) ; $b$ = 2,3~\% (1,7 à 3) ; $\lambda$ = 55~\% (40 à 70) & Recette de la fin de
l'effacement & Proportionnelle à chacun ; $b$ et $\lambda$ ensemble : recette de l'an 10 de 2,3 à 5,6~Md€. Pas de
statistique sur la composition des successions depuis 2006 \\
Ventes des héritiers : $a$ = 30~\%, puis $h$ = 4~\% & Profil de la recette & Les héritiers garderont plus leurs
titres : le jeu bas ramène l'an 10 de 3,6 à 2,1~Md€ ; les donateurs n'attendront plus le décès pour vendre (non
chiffré) \\
90~\% de la recette payée par le 1~\% & Charge du 1~\% & À 75~\%, 2,1~Md€ au lieu de 2,5. Le dernier centile paie
72~\% de l'impôt hors barème de 2025 (78~\% en 2024) \\
Impôt personnel du 1~\% sur le capital : 16~Md€ (21 en 2024, année gonflée par la CDHR) & « 20~\% » & Assiette du
centile déduite de son impôt hors barème (47~Md€ en 2025, 63 en 2024), plus 65~\% de la CEHR et la CDHR \\
Stock de reports de 144~Md€ (127 à 173), transmis en 25 ans (15 à 35), imposé à 30~\% & Report & Taux non
tranché ; abattements sur les entrées de 2013-2017 \\
CEHR payée à 50 à 80~\% sur le capital & Majoration (0,7 à 1,1~Md€ ; 1,0 à 1,6 en 2024) & Inconnue \\
Constitutionnalité (plus-values acquises avant la réforme ; effet dès l'annonce) & Transition & Non expertisée \\
\bottomrule
\end{tabularx}
\end{table}

\begin{table}[!htbp]\centering\scriptsize
\caption{Chiffres non vérifiés sur une source primaire}\label{tab:pfu_a_nonverif}
\begin{tabularx}{\linewidth}{@{}>{\raggedright\arraybackslash}p{4.4cm}>{\raggedright\arraybackslash}p{3.3cm}>{\raggedright\arraybackslash}X@{}}
\toprule
Chiffre ou fait & Où il joue & Statut \\
\midrule
Revenus de 2025 (impôt hors barème, CEHR, CDHR, revenus déclarés) & Situation & DGFiP n°~53, provisoire ; assiette
par nature de revenu estimée par l'auteur \\
Taux total de 46~\% à l'entrée du top 0,1~\%, 26~\% pour le top 0,0002~\% (2016) & Diagnostic & IPP, note n°~92,
non lue (via le rapport n°~3056) \\
84 foyers, 23,5~Md€ mis en report en 2023 & Diagnostic & Via le rapport n°~3056 ; incohérent avec les 18,8~Md€
du top 0,01~\% \\
\emph{Exit tax} dégrevée après deux ou cinq ans, quinze avant 2019 & Transition & Sources secondaires \\
Droits de succession sur les entreprises différés avec intérêt & Intérêt sur l'impôt différé & Sources
secondaires ; taux non vérifié \\
Rôle de l'évaluation au 31 décembre 1976 dans l'abandon américain & Clause de droits acquis & Non vérifié \\
Régime allemand hors participations de 1~\% ; règles australiennes ; PIB américain & Exemples ; repères &
Non vérifiés sur les textes \\
CSG déductible de 6,8~\% ; barème des revenus de 2024 ; IS de 15~\% jusqu'à 42\,500~€ & Cas types & Non
revérifiés pour 2026 ; effets de second ordre \\
\bottomrule
\end{tabularx}
\end{table}
